%-------------------------------------------------------------------------------
%	SECTION TITLE
%-------------------------------------------------------------------------------
\cvsection{Open-Source Projects}


%-------------------------------------------------------------------------------
%	CONTENT
%-------------------------------------------------------------------------------
\begin{cventries}

%---------------------------------------------------------
  \cventry
    {Initiator and Maintainer} % Role
    {aigverse} % Title
    {\href{https://github.com/marcelwa/aigverse}{https://github.com/marcelwa/aigverse}} % Location
    {2024 -- Present} % Date(s)
    {
      \begin{cvitems} % Description(s)
        \item {A Python library for working with logic networks, synthesis, and optimization.}
        \item {High-performance C++ backend and interoperability with other logic synthesis tools.}
        \item {Integration with popular data science and machine learning libraries.}
      \end{cvitems}
    }

%---------------------------------------------------------
  \cventry
    {Initiator and Maintainer} % Role
    {fiction} % Title
    {\href{https://github.com/cda-tum/fiction}{https://github.com/cda-tum/fiction}} % Location
    {2018 -- Present} % Date(s)
    {
      \begin{cvitems} % Description(s)
        \item {Framework for Design Automation of Field-coupled Nanotechnologies.}
        \item {Open-source project for researchers and developers.}
        \item {Honored with the ISVLSI 2019 Best Research Demo Award.}
      \end{cvitems}
    }

%---------------------------------------------------------
  \cventry
    {Contributor} % Role
    {Munich Quantum Toolkit (MQT) Core} % Title
    {\href{https://github.com/munich-quantum-toolkit/core}{https://github.com/munich-quantum-toolkit/core}} % Location
    {2025 -- Present} % Date(s)
    {
      \begin{cvitems} % Description(s)
        \item {Foundational C++/Python library of the Munich Quantum Toolkit, providing the shared data structures and algorithms (compilation, simulation, verification) that the toolkit's other quantum-software components build on.}
        \item {Contributing improvements and bug fixes as part of ongoing quantum software engineering work at MQSC.}
      \end{cvitems}
    }

%---------------------------------------------------------
  \cventry
    {Contributor and co-Maintainer} % Role
    {QDMI-on-IQM} % Title
    {\href{https://github.com/iqm-finland/QDMI-on-IQM}{https://github.com/iqm-finland/QDMI-on-IQM}} % Location
    {2025 -- Present} % Date(s)
    {
      \begin{cvitems} % Description(s)
        \item {Implementation of the Quantum Device Management Interface (QDMI) for IQM's quantum computing hardware, the standard interface connecting HPC software stacks to quantum backends.}
        \item {Enables the HPCQC integration case study demonstrated on real IQM hardware~\cite{arXiv_hpcqc_2026}.}
        \item {Co-maintain the repository: reviewing contributions and keeping the integration current with both QDMI and IQM's hardware interfaces.}
      \end{cvitems}
    }

%---------------------------------------------------------
\end{cventries}
